\documentclass[10pt,a4paper]{amsart}
\usepackage[margin=19mm]{geometry}
\usepackage{amsmath,amssymb,mathtools}
\usepackage[scr=rsfs,cal=euler]{mathalfa}
\usepackage{xfrac}
\usepackage[unicode,hidelinks,bookmarks=false]{hyperref}
\newcommand{\ie}{i.e.\ }
\newcommand{\cf}{cf.\ }
\newcommand{\resp}{respectively\ }
\newcommand{\Subsection}{\subsection}
\providecommand{\nobreakdash}{\nobreak\hbox{-}}
\setlength{\emergencystretch}{2em}
\multlinegap0pt
\usepackage{bm}
\usepackage{bbm}
\usepackage{dsfont}
\usepackage{xspace}
\usepackage{cancel}
\usepackage{mathtools}
\usepackage{amsthm}
\makeatletter
\@ifundefined{theorem}{\newtheorem{theorem}{Theorem}}{}
\@ifundefined{lemma}{\newtheorem{lemma}[theorem]{Lemma}}{}
\makeatother
\newcommand\johnord[1]{\mathcal{J}^{<}(#1)}
\newcommand\qup[1]{\mathbb{Q}_\nearrow^{#1}}
\newcommand{\Sym}{\operatorname{Sym}}
\newcommand{\aut}{\operatorname{Aut}}
\newcommand{\shuff}{\oplus}
\title[Taking model-complete cores]{Taking model-complete cores}
\author{Manuel Bodirsky, Bertalan Bodor, Paolo Marimon}
\date{Source: arXiv:2512.21278v2 (CC BY 4.0)}
\begin{document}
\maketitle

\noindent\emph{Source and licence.} Taking model-complete cores. Version arXiv:2512.21278v2 (CC BY 4.0), distributed under the Creative Commons Attribution 4.0 International licence, as recorded in the arXiv licence metadata for this exact version and shown on the abstract page https://arxiv.org/abs/2512.21278v2. Full rights notice: licenses/arxiv-2512.21278-CC-BY-4.0.txt.\par\smallskip

\noindent\textit{Editorial scope.} Counted unit: the Lemma whose source label is \texttt{lexicographic} in the article (source0.tex lines 1993--1996, source label \texttt{lexicographic}), delivered together with the article's own complete original proof of it (source0.tex lines 1998--2046, one \texttt{proof} environment of 49 lines). No mathematics is written, repaired or re-derived: statement and proof are the article's own text line for line. Required context retained uncounted: most\_sign (lines 1961--1964). Every definition, display and label of those lines is retained unchanged. Omitted from this package and not redistributed: 1--1960, 1965--1992, 1997--1997, 2047--2594. Nothing inside a retained excerpt is omitted or altered: every retained body is delivered exactly as the article's lines are written, so no omission receipt of any kind accompanies this package. Citations: the retained excerpts carry no citation command at all, so no bibliography entry of the article is transcribed and no reference is redistributed. Reference adaptation: none was needed and none was made; every reference inside the delivered unit resolves inside it, so no unresolved reference or citation is produced. Printed slips kept literal: no printed word, symbol, hypothesis or number of the counted statement or of its proof was altered. Layout and class: the article's own class is \texttt{amsart} with packages including graphicx, mathtools, amssymb, geometry, tikz, verbatim, hyperref, babel, wasysym, enumitem, todonotes, colortbl, etoolbox; the wrapper is the lane's pinned \texttt{amsart} editorial wrapper at 10pt on A4, which restates the theorem-like environments the retained text needs and adds the article's own macro definitions that the retained text needs (\texttt{\textbackslash Sym}, \texttt{\textbackslash aut}, \texttt{\textbackslash johnord}, \texttt{\textbackslash qup}, \texttt{\textbackslash shuff}), with extra wrapper packages (bm, bbm, dsfont, xspace, cancel, mathtools). The delivered numbering is the unit's own. Assets: no retained span contains a figure environment, a graphics-inclusion command, a PDF-page inclusion, a table or a TikZ picture; no image file is copied into this package, the package declares no asset dependency and no image file is redistributed with it. Licence: version arXiv:2512.21278v2 of this article is distributed under the Creative Commons Attribution 4.0 International licence, as recorded in the arXiv licence metadata for this exact version and shown on the abstract page https://arxiv.org/abs/2512.21278v2. A full rights notice is delivered at licenses/arxiv-2512.21278-CC-BY-4.0.txt. Characters: every retained excerpt is pure ASCII, so the delivered engine, XeLaTeX with its default Latin Modern text font, reports no missing character.
\begin{lemma}\label{most_sign}
	Let $d\in \mathbb{N}$, let $I_1<I_2<\dots<I_d$ be some non-degenerate intervals in $\mathbb{Q}$,
    and let $\prec$ be an order which is first-order definable in $(\qup{d};<_1,\dots,<_d)$. Then there are $k\in [d]$ and $R\in \{\prec,\succ\}$ such that for all $\bar{a},\bar{b}\in \mathcal{I}\coloneqq \prod_{i=1}^dI_i$ the relation $a_k<b_k$ implies $\bar{a}R\bar{b}$.
\end{lemma}

\begin{lemma}\label{lexicographic}
	Let $d\in \mathbb{N}$ and let $\prec$ be an order which is first-order definable in $\johnord{d}$. Then there are  $R_1,\dots,R_d\in \{<,>\}$ and $\sigma \in \Sym([d])$ such that for all $\bar{a},\bar{b}\in \qup{d}$ we have $\bar{a}\prec \bar{b}$ if and only if there exists some $j\leq d$ such that 
    $\{\sigma(1),\dots,\sigma({j-1})\} \subseteq (\bar{a}=\bar{b})$, and $a_{\sigma(j)}R_jb_{\sigma(j)}$.
\end{lemma}

\begin{proof}
	We show by induction on $m \in {\mathbb N}$, $m \leq d$, that we can choose some $R_1,\dots,R_m \in \{<,>\}$ and distinct $\sigma(1),\dots,\sigma(m) \in [d]$ such that for all $\bar{a},\bar{b}\in \qup{d}$, whenever there exists some $j\leq m$ such that $(\bar{a}=\bar{b})\supseteq \{\sigma(1),\dots,\sigma(j-1)\}$ and $a_{\sigma(j)} R_j b_{\sigma(j)}$, then $\bar{a}\prec \bar{b}$. $(\dag_m)$.
	
	For $m=d$ we obtain the statement of the lemma.

	We prove the statement by induction on $m$. For $m=0$ there is nothing to prove.
	
	For the induction hypothesis let us assume we already chose some $R_1,\dots,R_m$ and $\sigma(1),\dots,\sigma(m)$ such that $(\dag_m)$ holds. We introduce the following notation for ``shuffling'' tuples. 
    For $\bar{u}\in \mathbb{Q}^{d-m}$ and $\bar{c}\in \mathbb{Q}^m$ we write $\bar{u}\shuff_{\sigma} \bar{c}$ for the tuple whose $\sigma(i)$-th coordinate is $c_i$ for all $i \in [m]$, and such that the tuple formed by listing all the other coordinates in order is $\bar{u}$. Let us fix $\bar{c}\in \mathbb{Q}^m$ such that $$D\coloneqq \{\bar{u}: \bar{u}\shuff_{\sigma} \bar{c}\in \qup{d}\}\neq \emptyset.$$

	Note that in this case $\{\bar{u}\shuff_{\sigma} \bar{c}: \bar{u}\in \mathbb{Q}^{d-m}\}\cap \qup{d}$ is a $\prec$-interval by the induction hypothesis. Let us choose some non-degenerate intervals $I_1<\dots<I_{d-m}$ of $\mathbb{Q}$ such that $\mathcal{I}\coloneqq \prod_{i=1}^{d-m}I_i\subseteq D$. Note that for $\bar{u},\bar{v}\in D$ the $((<_{ij}))_{i,j\in [d]}$-type of the pair $(\bar{u}\shuff_{\sigma} \bar{c}, \bar{v}\shuff_{\sigma} \bar{c})$ is uniquely determined by the $((<_i))_{i\in [d]}$-type of the pair $(\bar{u},\bar{v})$. This implies that the order $$\prec^*\coloneqq \{(\bar{u},\bar{v}): \bar{u}\shuff_{\sigma} \bar{c}\prec \bar{v}\shuff_{\sigma} \bar{c})\}$$ is definable in $(\mathcal{I};<_1,\dots,<_d)$. Therefore, by Lemma~\ref{most_sign} we obtain that there exists some $k$ and some $R\in \{<,>\}$ such that $u_kRv_k$ implies $\bar{u}\prec^*\bar{v}$ for all $\bar{u},\bar{v}\in \mathcal{I}$. Since all sequences of intervals $(I_1,\dots,I_{d-m})$ as above can be mapped to each other by an automorphism of $\aut(\mathbb{Q})$ fixing $\bar{c}$, it follows that in fact the {relation $R$ and the} value $k$ does not depend on our particular choice of the intervals $I_i$.
	
	Now we show that in fact for all $\bar{u},\bar{v}\in D$ the relation $u_kRv_k$ implies $\bar{u}\prec^*\bar{v}$. For simplifying the notation we assume that $R$ is $<$; the proof of the other case is analogous. So suppose that $u_k < v_k$. We define the tuples $\bar{u}^0,\dots,\bar{u}^d,\bar{v}^0,\dots,\bar{v}^d$
    %\todo{BB: I think $d+1$ tuples are enough here, I'm not sure why I did this with infinitely many.} 
    as follows. Let $J$
    %\todo{MB: we already used $I$ above. BB: Changed it to $J$ but I also changed $I$ back to $\mathcal{I}$ (since it's not an interval but a product).} 
    be an inclusion-wise minimal non-degenerate interval {such that any of its existing endpoints belongs to $\bar{c}$.}
% whose endpoints \pau{(if they exist)}
    %(\green{whichever exist})
   % are elements of $\bar{c}$. 
   Let $u_i<u_{i+1}<\dots<u_j$ be the elements of $\bar{u}$ which are contained in $J$. Note that in this case $v_i<v_{i+1}<\dots<v_j$ are also exactly the elements of $\bar{v}$ in $J$. Let us pick some elements $w_i<\dots<w_j$ {in $J$} such that
	
\begin{itemize}
\item if $\ell<k$ then $w_{\ell}<\min(u_i,v_i)$,
\item if $\ell>k$ then $w_{\ell}>\max(u_j,v_j)$,
\item $u_k<w_k<v_k$.
\end{itemize}
For $\ell\in \{i,\dots,j\} \setminus \{k\}$ we define
\[
u^n_{\ell}\coloneqq 
\begin{cases*}
w_{\ell} & if $\ell<\min(k,i+n)$ or $\ell>\max(k,j-n)$ \\
u_{\ell} & otherwise.
\end{cases*}
\]
	The element $v^n_{\ell}$ is defined analogously. Finally, we define the elements $u^n_k$ and $v^n_k$ in such a way that
\[
u_k=u^0_k<u^1_k<\dots<u^d_k<u_{k+1},w_k,
\]
	and
\[
v_k=v^0_k>v^1_k>\dots>v^d_k>v_{k-1},w_k
\]
	{where $u_{d-m+1}$ and $v_0$ are considered to be $\infty$ and $-\infty$, respectively.}

	Note that by definition all tuples $\bar{u}^n$, $\bar{v}^n$, and $\bar{w}$ are in $D$, because they all have the same type over $\bar{c}$. Moreover, it follows from our choice that for all $n\in \mathbb{N}$ the intervals $[[u^n_{\ell},u^{n+1}_{\ell}]]: \ell\in [d-m]$ are pairwise disjoint. By making these intervals slightly bigger {we can assume that} we are in the setting of Lemma~\ref{most_sign}. Since ${u^n_k<u^{n+1}_k}$, we have that $\bar{u}^n\prec^*\bar{u}^{n+1}$. A similar argument shows that $\bar{u}^d\prec^*\bar{w}$. Therefore, we get $\bar{u}\prec^*\bar{w}$. By repeating the argument for the $v$-tuples we obtain $\bar{w}\prec^*\bar{v}$. Therefore, $\bar{u}\prec^*\bar{v}$, as we claimed.

	Now we can finish the induction step. We pick $\sigma({m+1})$ to be the position of the $k$-th coordinate after applying the shuffling map $\bar{u}\mapsto \bar{u}\shuff_{\sigma}\bar{c}$, and we put $R_{m+1}\coloneqq R$. Then clearly $\sigma(m+1) \not\in \{\sigma(1),\dots,\sigma(m)\}$. Now let $\bar{a},\bar{b}\in \qup{d}$ be arbitrary. If $(\bar{a}=\bar{b})\not\supseteq \{\sigma(1),\dots,\sigma(m)\}$ then we are done by the induction hypothesis. Otherwise, there exist some $\bar{u},\bar{v}\in \qup{d-m}$ and $\bar{c}'\in \qup{m}$ such that $\bar{a}=\bar{u}\shuff \bar{c}'$ and $\bar{b}=\bar{v}\shuff \bar{c}'$. Let us assume that $a_{\sigma(m+1)}Rb_{\sigma(m+1)}$. We have to show that $\bar{a}\prec \bar{b}$. This is enough to show in the case where  $\bar{c}'=\bar{c}$. In this case $u_k<v_k$, and therefore by our claim above $\bar{u}\prec^*\bar{v}$, that is, $\bar{a}=\bar{u}\shuff \bar{c}\prec \bar{v}\shuff \bar{c}=\bar{b}$.
\end{proof}

\end{document}
